\documentclass[11pt,a4paper]{article}
\usepackage[utf8]{inputenc}
\usepackage[margin=1in]{geometry}
\usepackage{fancyhdr}
\usepackage{graphicx}
\usepackage{amsmath}
\usepackage{amssymb}
\usepackage{array}
\usepackage{booktabs}
\usepackage{longtable}
\usepackage[most]{tcolorbox}
\usepackage{hyperref}
\usepackage{xcolor}

\pagestyle{fancy}
\lhead{Commodity Hedging \& Structuring Workbench}
\rhead{\today}
\cfoot{\thepage}

\newtcolorbox{concept}[1]{colback=black!4,colframe=black!60,
  fonttitle=\bfseries,title={#1},breakable}
\newtcolorbox{pitfall}[1][Pitfall]{colback=red!4,colframe=red!55!black,
  fonttitle=\bfseries,title={#1},breakable}

\title{Commodity Hedging \& Structuring Workbench\\
\large A Learning Reference and Validation Report}
\author{learner workbench (extension of the Citi MQA Forage simulation)}
\date{September 2026}

\newcommand{\usd}{\$}

\begin{document}

\maketitle

\section*{Executive summary}


\textbf{Validation verdict:} six of seven outcome checks PASS. The one
FLAG (historical 95\% VaR coverage of 38 breaches versus 25.2 expected
under the static historical estimator) is retained and discussed; the
filtered (EWMA-scaled) estimator remediated it, covering at 5.5\% of days
versus the 5\% expectation (ratio 1.10, inside the pre-declared band).

This report closes the validation phase of a practitioner-style commodity
hedging workbench built as an extension of the Citi MQA Forage simulation.
The workbench prices and hedges a roaster's coffee purchase program on a
frozen ICE Coffee~C (KC=F) curve: data gates, implied convenience-yield
term structure, a Schwartz--Smith two-factor fit, GARCH(1,1) working
volatility, a roaster futures-and-collar hedge program, a capital-protected
participation note, and a junior xVA layer (EE/EPE/PFE, unilateral CVA,
netting and collateral).

Phase~6 adds two things. First, an SR~26-2-aligned validation harness: a
model inventory, a fail-closed no-look-ahead gate on calibration windows,
and a rolling-window outcomes program across every model. Second, the
decision analysis: the roaster's outcome unhedged versus futures-only
versus a zero-cost collar, with margin-funding peaks per strategy.

\textbf{Headline decision numbers} (1.2M lb annual purchases, 2.67 average
contracts, front price 324.25~¢/lb): the full futures hedge locks the
annual purchase cost at \usd{}3,891,000 with a 95th-percentile margin peak
of \usd{}103,797; the zero-cost collar keeps the cost locked between the
strikes with a materially lower margin peak (\usd{}44,250 at the 95th
percentile) at the price of re-exposure beyond the strikes; the unhedged
roaster faces a 95th-percentile cost of \usd{}5,288,885 and a simulated
worst case of \usd{}10.3M.


All regulation content (SR 26-2, IFRS 9, EMIR 3, MiFID II/PRIIPs) is
\emph{mapped as methodology}, never asserted as compliance. All numbers in
this report are regenerated in one run by \texttt{report/make\_numbers.py}
from the frozen data (SHA-256 manifests); the hand-calculated checks of
Section~\ref{sec:handcalc} are reproducible without reading code.

\clearpage
\part{Concepts and methodology: how the workbench works}

Part~I explains every concept the workbench uses, in the order a
practitioner meets them: the business problem, the data, the curve, the
price model, the hedge, the note, the counterparty risk, and the regulation
that shapes each. Every concept follows the same pattern: what it is, why a
desk needs it, the formula, a worked number from the frozen data, and the
interpretation. Part~II then reports what was validated and what the numbers
decided.

\section{The business problem}

A coffee roaster buys green coffee continuously. Our reference program is
100{,}000~lb per month for twelve months (1.2M~lb total), priced off the ICE
Coffee~C futures contract (37{,}500~lb per contract). Coffee can move tens of
percent in a year; the roaster's retail margin cannot. Left unhedged, the
annual purchase cost of this program has a simulated 95th-percentile of
\usd{}5.29M and a worst case of \usd{}10.3M against a forward-implied
\usd{}3.89M --- that gap is the risk the workbench prices, hedges, and
explains.

\begin{concept}{Forward and futures contracts}
A \emph{forward} is a bilateral agreement to exchange at a fixed price on a
fixed date; a \emph{futures} contract is the exchange-traded equivalent,
standardized in size (\usd{}37{,}500 per 1~cent move per coffee contract),
expiration months, and quality. Two differences matter here. First,
futures are cleared through a central counterparty, so the roaster's
credit exposure to any single dealer is replaced by margin obligations to
the clearing house. Second, futures are marked to market daily: gains and
losses settle in cash every day (variation margin). The daily settlement
is why margin funding --- not just the final price --- becomes a risk to
manage.
\end{concept}

\begin{concept}{Basis}
The roaster's physical purchases settle at an origin differential over the
exchange~C price, not at the exchange price itself. \emph{Basis} is the gap
between the local physical price and the futures reference. A futures hedge
removes the exchange-price risk; the basis remains in the roaster's P\&L.
The workbench evaluates every hedge on the exchange series and says so:
basis is not designated in the hedge relationship and not modelled
(Section~\ref{sec:limitations}).
\end{concept}

\begin{concept}{Backwardation and contango}
A futures \emph{curve} is a set of prices by expiry month. \emph{Backwardation}
means deferred months trade below the front; \emph{contango} is the reverse.
On the frozen curve (as-of 2026-09-04) every consecutive spread is
backwardated: the front (Sep~2026) is 324.25~\textcent/lb and each later month
sits lower. A backwardated curve is the commodity analogue of a stock paying
a large dividend: something valuable accrues to holders of the physical
good.
\end{concept}

\begin{concept}{Convenience yield}
That ``something'' is the \emph{convenience yield}: the benefit of holding
physical inventory (ability to ship now, keep roasting lines fed, exploit
local scarcity) beyond its financing and storage cost. When near coffee is
scarce, users bid up the front of the curve and the implied convenience
yield rises. Chapter~\ref{ch:carry} extracts it from the curve; the frozen
curve implies 40.8\%/yr at the front, decaying to 6.3\%/yr on the back ---
the market is paying heavily for immediate availability.
\end{concept}

Why hedge at all: the unhedged cost distribution above is one tail away
from a bad year. The rest of Part~I builds, piece by piece, the models that
turn market data into that distribution and into hedge structures that
change it.

\section{Data engineering as a model-risk control}
\label{ch:data}

Every number in this workbench traces to a small set of frozen files. That
is deliberate: in model-risk terms, uncontrolled input data is itself a
model risk, and the cheapest control is a reproducible, checksummed
snapshot.

\begin{concept}{Data freeze with SHA-256 manifests}
The downloader writes each series to \texttt{data/frozen/}, computes a
SHA-256 hash per file, and records it in a JSON manifest. Every later use
of the data re-verifies the hash before reading. A changed file is not
silently accepted; the pipeline stops. This is what makes ``every number
regenerated in one run'' a verifiable claim rather than a hope: the
manifests are committed, so any reader can confirm the exact bytes behind
the results.
\end{concept}

\begin{concept}{No-look-ahead gate}
A \emph{look-ahead error} occurs when a calibration window accidentally
sees data after its valuation as-of date; it silently inflates every
backtest. The workbench makes the failure impossible to miss: each
calibration call passes through a gate that truncates the series to the
as-of window and raises a \texttt{LookaheadError} if any observation past
the as-of date is requested. The gate is itself tested by feeding it a
poisoned window and asserting it refuses.
\end{concept}

\begin{concept}{Fail-closed design with a pre-declared fallback}
The universe logic is \emph{fail-closed}: coffee gate failures switch to
the pre-declared gold fallback; a gold failure stops everything and
nothing downstream is produced. The opposite (fail-open) design, returning
partial or stale data on failure, is how mismatched numbers enter
production systems unnoticed. Declaring the fallback \emph{before} it is
needed prevents the subtle failure where the fallback is chosen because it
gives the preferred answer.
\end{concept}

The sources are stated plainly: Yahoo Finance daily bars (an unofficial
source, accepted and documented) for the ICE contract chain, and the SOFR
rate from FRED \cite{fredsofr}. The manifests, gate reports, and executed
notebooks are committed; the price bars are not (exchange-data licensing,
Section~\ref{sec:scope}) and regenerate locally with one command.

\section{The forward curve and cost of carry}
\label{ch:carry}

The curve is where commodity desks start. This chapter derives the one
identity that turns two futures prices into an implied yield, applies it to
the frozen coffee curve, and reads what it says about the physical market.

\begin{concept}{The cost-of-carry identity, derived}
Suppose you buy coffee now at spot $S$, pay financing at rate $r$ and
storage at rate $d$, and hold one unit for a period $\tau$; a futures sale
at $F$ locks the outcome. No-arbitrage pins the two choices together:
buying and carrying must cost the same as buying forward, so
$F = S\,e^{(r+d-y)\tau}$, where $y$ is the annual convenience yield
introduced in Chapter~1. The convenience yield enters with a minus sign:
holding the physical good earns the convenience benefit, so the forward
only has to compensate financing plus storage \emph{minus} that benefit.
Now take the ratio of two adjacent futures prices $F_1$ (near) and $F_2$
(next expiry, $\tau$ apart). Spot, $S$, cancels, as does everything not
known between the two dates:
\[
  \frac{\ln F_2 - \ln F_1}{\tau} \;=\; r + d - y
  \qquad\Longrightarrow\qquad
  y \;=\; r + d - \frac{\ln(F_2/F_1)}{\tau}.
\]
Two adjacent futures prices and a financing rate are enough to extract the
market's own implied convenience yield.
\end{concept}

\begin{concept}{Worked on the frozen curve}
As of 2026-09-04 the coffee chain carries 8 contract months plus the
continuous front (KC=F). With SOFR frozen at 3.66\% and storage assumed
zero (see the caveat below), applying the identity spread by spread gives
an implied convenience yield of 40.8\%/yr at the front, decaying to
6.3\%/yr on the back, with the curve backwardated at every consecutive
spread. The package computes this as \texttt{carry.implied\_yield}; the
curve-stability outcomes check re-runs it at six historical as-of dates
(front yield 14.8--40.8\%/yr; backwardated at 4 of 6, with two historical
contango episodes preserved rather than tuned away).
\end{concept}

\emph{Interpretation.} A 40\%/yr implied convenience yield says the market
prices immediate availability as scarce: hedgers who need physical coffee
now, or who fear not having it, bid the front up relative to deferred
months. The decay from front to back is normal --- scarcity is a
short-dated phenomenon --- and the backwardated shape is the term
structure signature of that scarcity. For the roaster program this matters
directly: rolling a long position through a backwardated curve
(\emph{rolling up} the curve) tends to buy high and sell low each cycle, a
recurring carry cost the hedge program must budget for.

\begin{pitfall}[Storage is assumed, not observed]
The identity reports $y$ \emph{gross of storage}: separating convenience
yield from storage outlay $d$ needs a storage-cost assumption the data
does not carry, so the workbench fixes $d=0$ and states it. Reported
``convenience yield'' is therefore an upper bound on the true yield by the
size of the true storage rate. The honest fix is not a secret calibration
but a stated assumption.
\end{pitfall}

\begin{pitfall}[An unofficial source can distort a yield]
Yahoo prices are unofficial; a stale or inconsistent quote between
contract months would show up as an implausible implied yield. This is why
the curve-stability check exists: it re-derives the yield at six as-of
dates and asserts the front yield stays inside a pre-declared plausibility
band. A data problem announces itself as a curve problem.
\end{pitfall}

\section{Price dynamics: level and volatility}
\label{ch:dynamics}

Three later chapters need a price process: the exposure engine simulates
future MtM, the Monte Carlo prices option payoffs, and the hedge analysis
stress-tests funding peaks. This chapter covers the two models the
workbench fits --- a two-factor level model and a volatility model --- and
the discipline of choosing between them with pre-declared criteria.

\subsection{Level: the Schwartz--Smith two-factor model}

\begin{concept}{Two factors, two horizons}
Commodity prices mix two behaviors: short-term mean reversion (weather,
supply shocks, inventory swings wash out) and long-term drift in the
equilibrium price level (cost inflation, demand cycles). Schwartz and
Smith \cite{schwartzsmith00} model the log price as the sum of both:
$S_t = \exp(X_t + \chi_t)$, where $X_t$ is the long-term factor (a
geometric walk with equilibrium level $\xi$ and reversion speed $\kappa$)
and $\chi_t$ is the short-term factor (an Ornstein--Uhlenbeck deviation
that decays toward zero at speed $\kappa$). The mean-reversion speed has a
natural unit: the \emph{half-life} $\ln 2 / \kappa$, the time for half of
a deviation to disappear.
\end{concept}

\begin{concept}{Fit and interpretation on the frozen curve}
The reduced-form fit on the frozen curve gives mean-reversion speed
$\kappa = 5.97$/yr, half-life $\ln 2 / 5.97 \approx 0.116$~yr
($\approx 1.4$ months), and log equilibrium level $\xi = 5.673$, i.e.\
$e^{\xi} \approx 291$~\textcent/lb. The fit RMSE sits far inside the
pre-declared tolerance. Interpretation: the market believes coffee
reverts to roughly 291~\textcent/lb and that price surprises are absorbed
in weeks, not months --- while the current front (324.25) stands above the
equilibrium, consistent with the scarcity reading of the 40\% front
convenience yield.
\end{concept}

\begin{pitfall}[Reduced form, single snapshot]
The fit uses one curve snapshot, which identifies the level parameters
$(\kappa, \xi)$ but not the volatility split between the two factors; a
full stochastic fit needs Kalman maximum-likelihood estimation on the
whole panel of maturities. The workbench flags this rather than hiding
it: the fit is used for level interpretation, and every premium in
Part~II carries a GARCH volatility estimate on top of it.
\end{pitfall}

\subsection{Volatility: EWMA and GARCH(1,1)}

\begin{concept}{EWMA --- the RiskMetrics estimator}
The exponentially weighted moving average \cite{riskmetrics96} updates
yesterday's variance with today's squared return,
$\sigma_t^2 = \lambda\,\sigma_{t-1}^2 + (1-\lambda)\,r_{t-1}^2$, with
decay $\lambda = 0.94$ on daily data. It is simple, always defined, and
reacts fast to volatility clustering --- but it has no long-run level: it
is a GARCH(1,1) with persistence forced to 1, so it can drift far from
anything sustainable. On the frozen returns the EWMA reads 47.2\%/yr.
\end{concept}

\begin{concept}{GARCH(1,1) --- the workhorse}
The GARCH(1,1) of Bollerslev \cite{bollerslev86} adds mean reversion:
$\sigma_t^2 = \omega + \alpha\,r_{t-1}^2 + \beta\,\sigma_{t-1}^2$. The
long-run variance is $\omega/(1-\alpha-\beta)$, and $\alpha+\beta$ is the
\emph{persistence} --- how slowly variance decays back to it. On the
frozen data: persistence $0.802$, long-run level 38.5\%/yr, last
conditional estimate 39.2\%/yr. The current level sits almost exactly at
its own long-run mean --- the market is in a normal-volatility regime.
\end{concept}

\begin{concept}{Choosing with a pre-declared criterion}
Model choice is an outcome, not a preference: the workbench declares
before looking that GARCH is adopted only if its rolling one-step forecast
RMSE beats the EWMA reference with ratio $\leq 1.1$. On four rolling
windows the RMSEs are 4.01 vs 6.80\%/yr (ratio 0.59) --- GARCH wins
clearly and is adopted as the working volatility. The EWMA stays as the
benchmark model in the inventory, which is exactly how an
SR~26-2-style outcomes analysis should treat a model selection
(Section~\ref{sec:regulatory}).
\end{concept}

\begin{pitfall}[Historical volatility is not implied volatility]
Every premium and VaR below inherits this one number from frozen history;
parameter risk flows through. No free implied-volatility surface exists
for coffee options, so there is no market check on the level --- stated
as a limitation in Section~\ref{sec:limitations}, not fixed by hoping.
\end{pitfall}

\section{Hedging the roaster program}
\label{ch:hedge}

With a curve, a carry reading, and a volatility model, the roaster can
actually hedge. This chapter builds the two candidate structures --- the
plain futures lock and the zero-cost collar --- and states exactly what
each one does to the cost distribution.

\subsection{The futures hedge and its plumbing}

The program is 1.2M~lb over 12 months, i.e.\ 32 contract-months
($1.2\mathrm{M}/37{,}500$), an average of $2.67$ contracts outstanding.
The hedge selects each month's contract with the \emph{in-or-after} rule
(the expiry month at or after the physical delivery date) and rolls each
contract 10 trading days before expiry, so no position is forced into
delivery.

\begin{concept}{Margin: the liquidity side of hedging}
Posting the hedge requires \emph{initial margin} (performance deposit,
\usd{}8k per contract here --- a stated learner assumption, not an
exchange feed) and daily \emph{variation margin}: every day the futures
lose, cash leaves. A perfectly correct hedge can still bankrupt its owner
in the middle if the price moves far enough before it pays off. This is
why the workbench reports \emph{margin-funding peaks} (the running
cumulative cash drain) alongside every cost number: the futures hedge
locks the price at \usd{}3{,}891{,}000 but carries a 95th-percentile
margin peak of \usd{}103{,}797 that must be funded throughout.
\end{concept}

\subsection{The zero-cost collar}

\begin{concept}{Construction}
Buy the futures hedge, buy a put struck at $K_p = 280$, and sell a call
struck at $K_c = 382.06$ with the premium equal to the put's (both
7.69~\textcent/lb; the strike $K_c$ is solved so the premiums net to zero
--- the residual gap is $1.3\times10^{-14}$ by construction). The put
pays $(K_p - f_T)^{+}$ dollars per lb at expiry; the short call costs
$(f_T - K_c)^{+}$. Per pound, the terminal cost is therefore
\[
  \text{cost}(f_T) \;=\;
  \begin{cases}
    f_0 - K_p + f_T & f_T < K_p \quad \text{(put monetized, cost floats but
      below the lock)}\\[2pt]
    f_0 & K_p \leq f_T \leq K_c \quad \text{(cost locked at the forward)}\\[2pt]
    f_0 - K_c + f_T & f_T > K_c \quad \text{(short call loses, cost floats)}
  \end{cases}
\]
and the hedge P\&L is capped between the strikes, which caps the
\emph{margin funding} at the put protection:
$(324.25 - 280)\times \tfrac{8}{3}\times 375 = \usd{}44{,}250$ at the
95th percentile, versus \usd{}103{,}797 for the futures lock.
\end{concept}

\begin{pitfall}[The strikes bound the hedge P\&L, not the buyer's cost]
``Locked between the strikes'' is easy to misread. The \emph{hedge
position} gains are capped; the \emph{buyer's total cost} still floats
1:1 with the market outside them: below 280 the monetized put drags the
cost back down (a crash keeps helping), above 382.06 the short call
bleeds (a spike keeps hurting). The collar is not a budget guarantee; it
is a trade that trades tail cost for liquidity --- p95 cost \usd{}4.60M
versus the futures' locked \usd{}3.89M, with margin peaks roughly halved.
\end{pitfall}

\emph{Choosing between them.} For a roaster whose binding constraint is
cash-flow certainty, pick the futures lock: same locked cost, full margin
exposure. For a margin-constrained buyer, the collar halves the funding
peak at the price of re-exposure beyond the strikes. Both structures are
priced in the decision section (Part~II) on the same scenario engine, so
the comparison is like-for-like. Hedge accounting for either is treated
in Section~\ref{sec:regulatory} (IFRS~9); the clearing-threshold
classification in the EMIR~3 part of the same section.

\section{The structured participation note}
\label{ch:note}

Structured notes are where curve, volatility, and discounting meet in one
instrument. The workbench builds the desk-side version of the product the
Citi Forage simulation gestures at: a note issued to an investor that
returns principal plus a share of the coffee rally.

\begin{concept}{Capital-protected participation payoff}
The note pays, per \usd{}1 of notional at maturity $T$:
\[
  1 + p\,\max\!\left(\frac{F_T}{F_0}-1,\;0\right),
  \qquad p = 0.6,
\]
a 9-month ATM participation on KC=F. The investor keeps the principal in
every scenario and keeps 60\% of any rally. This decomposes by design
into two ordinary instruments: a zero-coupon bond (the principal) plus
$p$ calls struck at the forward $F_0$.
\end{concept}

\begin{concept}{Black-76 pricing}
Options on futures are priced by Black's model \cite{black76}: the
premium is the discounted expectation under the risk-neutral measure in
which the futures is a martingale,
\[
  c \;=\; e^{-rT}\bigl(F\,N(d_1) - K\,N(d_2)\bigr), \qquad
  d_1 = \frac{\ln(F/K) + \sigma^2 T/2}{\sigma\sqrt{T}}, \quad
  d_2 = d_1 - \sigma\sqrt{T},
\]
with the GARCH working volatility $\sigma$ from
Chapter~\ref{ch:dynamics} and $r$ the frozen SOFR. The premiums of both
collar legs in Chapter~\ref{ch:hedge} come from the same formula, which
is why a zero-premium collar is solvable at all.
\end{concept}

\begin{concept}{Worked decomposition}
At $F_0 = 324.25$, 9-month tenor, SOFR 3.66\%: the discount bond is
$e^{-rT} = 0.97292$; the ATM Black-76 call is 0.13128, so the
participation piece is $0.6 \times 0.13128 = 0.07877$. Fair value:
$0.97292 + 0.07877 = 1.05169$ per unit notional. An investor paying
105.17\% of notional is financing 2.92 points of protection with 60\% of
the upside.
\end{concept}

\begin{concept}{Monte Carlo as an independent check}
The closed form assumes its own inputs; an implementation can still be
wrong. The workbench re-prices the note with an antithetic Monte Carlo
engine (100k paths, variance reduced by pairing each path's price with
its mirror image): $1.051913 \pm 0.000452$, against the closed form at
1.051689 --- a difference of $8.2\times10^{-5}$, under 2 standard errors.
Two engines, two code paths, one answer: that two-way match is the trust
mechanism, and it is asserted as an outcomes check in Part~II.
\end{concept}

\begin{pitfall}[One model, one premium]
Both the collar legs and the note inherit the GARCH point estimate; there
is no implied surface to cross-check against (no free listed coffee
options chain). The issuer-side greeks (delta, vega) are validated by
bump-and-reprice rather than market calibration. Parameter risk is
declared, not eliminated.
\end{pitfall}

Regulatory placement: a note like this sold to retail in the EU would
require a PRIIPs key information document
(Section~\ref{sec:regulatory}); here the note is priced, never offered.

\section{Counterparty risk and xVA}
\label{ch:xva}

Any issuer or dealer holding the other side of these positions carries
credit exposure: if the counterparty defaults while the trade is in our
favor, we recover less than full value. The junior xVA layer quantifies
that exposure and its price.

\begin{concept}{Exposure simulation: EE, EPE, PFE}
The engine simulates the futures price under the risk-neutral measure as
a lognormal martingale (the same assumption Black-76 prices), marks the
book's MtM on every path, and collects, per future date $t$, the
distribution of $\max(\text{MtM}, 0)$ --- the exposure only matters when
we are owed money. Three summaries:
\[
  \mathrm{EE}(t) = \mathbb{E}\bigl[\max(\text{MtM}(t),0)\bigr],
  \quad
  \mathrm{EPE} = \tfrac{1}{T}\!\int_0^T \mathrm{EE}(t)\,dt,
  \quad
  \mathrm{PFE}_{q}(t) = q\text{-quantile of } \max(\text{MtM}(t),0).
\]
EPE is the average exposure (what CVA prices); PFE is the tail exposure
(what collateral limits are set against). By construction
$\mathrm{EE}(t) \leq \mathrm{PFE}(t)$, and the suite asserts it.
On the 2.67-contract book: EPE \usd{}33{,}755; 1-year PFE \usd{}250{,}975.
\end{concept}

\begin{concept}{Unilateral CVA, bucketed}
Credit valuation adjustment is the discounted expected loss from
counterparty default, integrating exposure over time:
\[
  \mathrm{CVA} \;=\; \mathrm{LGD}
  \int_0^T e^{-rt}\,\mathrm{EE}(t)\; h\, e^{-ht}\,dt,
\]
with default hazard $h$ (flat 3\% here) and loss given default LGD = 60\%.
The integral is evaluated on discrete time buckets (the $h\,e^{-ht}$
term is the per-bucket default density). On this book: \usd{}583.79. The
engine is benchmarked by re-quoting a matched EE grid through QuantLib
\cite{quantlib} --- agreement to $1.9\times10^{-16}$ relative on that grid
(Section~\ref{sec:validation}); the notebook's matched pair on this same
engine (notebook 04) agrees to $4.8\times10^{-16}$.
\end{concept}

\begin{concept}{Netting}
Under a netting agreement, exposure to a set of positions is the net MtM
of the set, not the sum of gross exposures: offsetting positions cancel
before default law looks at them. Consequence: netted CVA
$\leq$ sum of standalone CVAs, asserted in the suite. For the roaster
program the netting set is the futures positions across contract months.
At desk scale this is why ISDA netting collapses multibillion gross
exposures to a fraction.
\end{concept}

\begin{concept}{Collateral}
Collateral --- a threshold (unsecured exposure the counterparty may build
before posting) plus independent amounts (posting required in advance)
--- attenuates the EE profile monotonically: every threshold unit cuts
the tail exactly once, everywhere. Collateral limits are set against
PFE, not EE: a limit sized on the average exposure is breached by
half the paths. Thresholds also transform CVA, and the transform is
directional (less unsecured exposure, less CVA) but not uniform.
\end{concept}

\begin{concept}{Wrong-way risk, read rather than modelled}
WWR is correlation between counterparty default and exposure. Here it is
benign by inspection: the roaster buys physical coffee and hedges long,
so its exposure to us grows when coffee rises --- while the roaster's own
business (it is a buyer of coffee) is stressed when coffee rises. That is
\emph{right}-way risk for us as the issuer of the note's short call. The
dangerous mirror is selling the participation note's upside to a producer
long coffee: exposure and counterparty distress would then rise
together. Modelled WWR needs a correlated default--price process; kept
qualitative here per scope, stated as such.
\end{concept}

\begin{pitfall}[Junior xVA only]
No DVA, FVA, or MVA; flat hazard and flat SOFR discounting. These
omissions matter at dealer scale (funding costs of the hedge itself,
own-default leg of the CVA). The layer is sized for a junior desk
review, and says so.
\end{pitfall}

\section{Regulatory context}
\label{sec:regulatory}

Four regimes shape what this workbench does. Each is treated the same
way: the problem the regime exists to solve, what it requires, and how
the workbench maps to it --- mapped as methodology, never asserted as
compliance, because a learner artifact has no regulatory status to
assert.

\subsection{SR 26-2: model risk management}

\begin{concept}{The problem and the guidance}
Models simplify reality; the gap between model and reality is model
risk --- the potential for loss or bad decisions when outputs are wrong
or misused. The US banking agencies first codified expectations in
SR~11-7 (2011) and issued the revised guidance SR~26-2 on 2026-04-17,
superseding SR~11-7 and the BSA/AML addendum SR~21-8 \cite{sr262}. The
revised guidance is explicitly risk-based and tailored: it is expected
to be most relevant to banking organizations above \usd{}30B in assets,
and it sets no enforceable standards.
\end{concept}

\begin{concept}{What SR 26-2 actually requires}
A \emph{model} is a quantitative method applying statistical, economic,
or financial theory to process input data into estimates --- simple
arithmetic and rule-based software are excluded. Requirements, by
lifecycle stage:
\begin{itemize}
\item \emph{Development}: clear statement of purpose, testing out-of-sample
  and out-of-time, sensitivity to assumptions.
\item \emph{Validation}: (i)~\emph{conceptual soundness} --- assess the
  design, assumptions, and judgments; (ii)~\emph{outcomes analysis} ---
  compare model outputs to real-world outcomes (backtests, benchmarks);
  (iii)~\emph{ongoing monitoring} --- detect performance deterioration.
\item \emph{Governance}: roles and independence (\emph{effective
  challenge} by qualified, objective reviewers), a comprehensive model
  inventory, documentation adequate to understand each model's risk.
\item \emph{Aggregate view}: risk assessed across models together ---
  common assumptions and shared data create correlated failures.
\item \emph{Vendor products}: third-party models are validated like
  in-house ones.
\end{itemize}
\end{concept}

\emph{Mapping.} The workbench reproduces the \emph{mechanics} on a small
scale: a 13-model inventory (Part~II, Table~\ref{tab:inventory}), a
rolling outcomes program across every model, benchmarking (GARCH vs
EWMA, MC vs closed form, CVA vs QuantLib), fail-closed data gates, and
hand-calculated checks in test docstrings as documentation. The posture
is methodological fidelity: practice the discipline, claim no status.

\subsection{IFRS 9: hedge accounting}

\begin{concept}{The problem}
Without hedge accounting, a hedging futures position marks to market in
P\&L \emph{today} while the hedged purchase hits earnings \emph{later};
earnings volatility appears purely from timing. Hedge accounting aligns
the two by parking the effective portion of the hedge's gain or loss in
other comprehensive income until the hedged transaction affects
earnings.
\end{concept}

\begin{concept}{Cash-flow hedge designation and effectiveness}
IFRS 9 \cite{ifrs9} permits designation of a forecast transaction as a
hedged item when it is \emph{highly probable} and presents a particular
risk. For the roaster program: highly probable forecast green-coffee
purchases (variable price), with coffee~C price risk designated --- the
basis between origin differentials and the exchange~C price is
\emph{not} designated and stays in P\&L. Effectiveness is principles-based:
there must be (i)~an economic relationship between hedged item and
instrument, (ii)~credit risk not dominating value changes, and
(iii)~a hedge ratio consistent with the risk-management objective ---
the retired IAS~39 80--125\% dollar-offset band is gone. The effective
portion accumulates in AOCI and is reclassified to P\&L in the period
the forecast purchases affect earnings. Variation margin posted to the
clearing house is a receivable/liability, not a hedge-accounting item.
\end{concept}

\emph{Mapping.} The program is assessed with
\texttt{effectiveness.assess} on the hedge P\&L pairs: economic
relationship and credit-dominance tests, boundary cases included;
forward-starting assessments run each reporting period. The designation
memo (conceptual, not compliance advice) documents the archetype: a
roaster/buyer patterned on how a branded coffee company describes its
commodity hedging in its annual report.

\subsection{EMIR 3: clearing thresholds}

\begin{concept}{The problem and the regime}
After 2009, G20 commitments pushed standardized derivatives to central
clearing and reporting; EMIR implements that in the EU. EMIR~3
(Regulation (EU) 2024/2987, in force 24 December 2024
\cite{emir3}) reshaped counterparty categorisation: financial
counterparties (FC) and non-financial counterparties (NFC+, NFC-)
are sorted by their derivative positions, and only counterparties
exceeding \emph{clearing thresholds} are forced to clear.
\end{concept}

\begin{concept}{What changed and where this program sits}
EMIR~3 moved the threshold calculation to \emph{uncleared OTC
positions only} (recognizing the benefit of clearing), computed at
entity level for NFCs, with the \emph{hedging exemption} retained:
contracts that objectively reduce commercial risks of the counterparty
or its group are excluded. ESMA's February~2026 final technical
standards set the uncleared commodity and emission-allowance threshold
at EUR~4 billion \cite{emir3rts}. This program's annual notional is
roughly EUR~3.6M ($\approx$ 0.009\% of the threshold): the clearing
obligation is nowhere near engaged, and the workbench's exchange-traded
futures are cleared by construction. The mapping is directional
context, not a classification claim.
\end{concept}

\subsection{MiFID II and PRIIPs: distribution and disclosure}

A half-section because it only frames the note's boundaries. MiFID~II
governs derivative trading and sets position limits for commodity
derivatives (with hedging exemptions); PRIIPs (Regulation (EU)
No~1286/2014, \cite{priips}) requires a key information document before
any structured product is offered to a retail investor. The
participation note in Chapter~\ref{ch:note} is priced two ways but
offered to nobody; the KID obligation maps to zero here. Recorded so
the boundary is explicit.

\clearpage
\part{Validation and decision: what was checked and what it decided}

\section{Scope and Honest Framing}\label{sec:scope}

This project extends the Citi MQA Forage simulation onto real frozen
market data. It is a learner workbench, not a production system: no Citi
affiliation is claimed, no internal methodology is reproduced, no
regulatory compliance is asserted, and nothing here is investment advice.
The data (Yahoo Finance daily bars, SOFR from FRED) is checksummed
and gated; exchange-data redistribution licensing forbids committing the
price bars, so only manifests, gate reports and executed notebooks are
committed and the bars regenerate locally through the gated freeze. Any
future cloud artifact is gated snapshots only.

\section{Program structure}

The workbench was built in five methodology phases and two review
phases. Part~I explains each phase's methods chapter by chapter; the
map below links phases to deliverables for navigation.

\begin{table}[h!]
\centering\small
\begin{tabular}{p{2.2cm}p{7.2cm}p{5.2cm}}
\toprule
\textbf{Phase} & \textbf{Method} & \textbf{Explained in} \\
\midrule
1. Data gates + MQA replication & SHA-256 manifest freeze, count/completeness gates, gold fallback, fail-closed; the Forage Task-3 methodology re-run on the frozen curve & Ch.~2 (\ref{ch:data}); outcomes in \ref{sec:validation} \\
2. Carry \& curve calibration & Cost-of-carry identity, Schwartz--Smith fit, GARCH(1,1) vs EWMA & Ch.~3 (\ref{ch:carry}), Ch.~4 (\ref{ch:dynamics}) \\
3. Roaster hedging program & Futures + zero-cost collar, contract selection and roll, margin funding & Ch.~5 (\ref{ch:hedge}) \\
4. Structured note & Black-76 closed form vs antithetic Monte Carlo, bump-and-reprice greeks & Ch.~6 (\ref{ch:note}) \\
5. Risk gates + junior xVA & EE/EPE/PFE engine, bucketed unilateral CVA, netting and collateral & Ch.~7 (\ref{ch:xva}) \\
6. Validation harness & SR 26-2-aligned inventory, no-look-ahead gate, rolling outcomes & \ref{sec:validation} \\
\bottomrule
\end{tabular}
\caption{Phase map. The decision analysis and limitations close Part~II.}
\end{table}

\section{Hand-Calculated Verification}\label{sec:handcalc}

Per the spec, an interviewer must be able to verify the numbers without
reading code. Each check below is a test in the suite with the arithmetic
worked out in its docstring.

\begin{table}[h!]
\centering\small
\begin{tabular}{p{4.2cm}p{7.6cm}p{3.4cm}}
\toprule
\textbf{Check} & \textbf{Hand calculation} & \textbf{Result} \\
\midrule
CVA, two buckets, constant EE \usd{}1M, hazard 10\%, LGD 60\%, $r=4$\% &
$dS_1 = 1-e^{-0.05} = 0.04877058$, $df_1 = e^{-0.02} = 0.98019867$;\newline
$dS_2 = e^{-0.05}-e^{-0.10} = 0.04639219$, $df_2 = e^{-0.04} = 0.96078944$;\newline
CVA $= 0.6 \times 10^6 \times (0.04780387 + 0.04457460)$ &
\usd{}55,427.08 (exact match) \\
\midrule
VaR coverage, two 5-point windows at 80\% &
W1 hist $[-5..-1]$: VaR $= |q_{0.2}| = 4.2$, next window 0 breaches;\newline
W2 hist $[-2..2]$: VaR $= |q_{0.2}| = 1.2$, next window ($-10$ first) 1 breach &
VaR $[4.2, 1.2]$, breaches $[0, 1]$ \\
\midrule
EWMA forecast, returns $[1,2,3]$\%, $\lambda = 0.94$ &
weights $w = (1, 0.94, 0.8836)$; weighted variance
$\frac{\sum w (x-\bar x)^2}{\sum w - \sum w^2 / \sum w} = 0.9993624$;
annualised $\sqrt{252 \times 0.9993624}$ &
15.869\%/yr \\
\midrule
Collar net cost, synthetic zero-gap collar (f0 320, put 280, call 350),
volume 100k lb &
unhedged $= 1000 f_T$;\newline futures locked at 320k;\newline
collar: $240$k below the put ($f_T - K_p + f_0$), 320k between, 350k at
$f_T = 380$ &
exact on all four scenarios \\
\midrule
Filtered VaR (95\%), 4-day price series 100/101/103/102,
\usd{}1000 per return-fraction &
VaR at $t$ uses only returns $\leq t{-}1$ and price($t{-}1$): EWMA var
($\lambda=0.94$, weights $0.06/1$) over $\{1\%, 1.9802\%\}$
$= 0.054381/0.11321 = 0.48039\,\%^2$ $\rightarrow$ sd $0.6931\,\%$/day;\newline
VaR $= 1.6449 \times 0.006931 \times 103 \times 1000$ &
1,174.26 (exact match) \\
\midrule
Collar margin cap, contracts $= 2$ (\usd{}750/cent) &
futures drawdowns 15k/90k; collar funding loss capped at
$(320-280)\times 750 = 30$k &
$[15k, 30k]$ capped \\
\bottomrule
\end{tabular}
\caption{Hand-calculated verification table (all reproduced exactly by the
test suite).}
\end{table}

\section{SR 26-2-Aligned Validation}\label{sec:validation}

The posture is methodological fidelity, not compliance: the harness
reproduces the \emph{mechanics} SR~26-2 (2026-04-17, superseding SR 11-7)
expects of model-risk discipline (inventory, conceptual soundness,
outcomes analysis, benchmarking) over a frozen dataset, with every
finding reproducible from the package.

\subsection{Model inventory}

\begin{table}[h!]\label{tab:inventory}
\centering\scriptsize
\begin{tabular}{p{2.7cm}p{2.5cm}p{3.4cm}p{3.2cm}p{2.7cm}p{1.5cm}}
\toprule
\textbf{Model} & \textbf{Inputs} & \textbf{Key assumptions} & \textbf{Benchmark} & \textbf{Outcomes check} & \textbf{Review} \\
\midrule
Data gates & frozen chain + manifests & count/completeness ceilings & SHA-256 checksums & fail-closed on injected gap & 2026-09, green \\
Implied carry & frozen chain, SOFR & expiry $\approx$ 15th; storage $d=0$ & cost-of-carry identity & curve stability at past as-ofs & 2026-09, green \\
Schwartz--Smith & curve snapshot & reduced form; $\kappa$ weakly identified & NLS RMSE, $\kappa>0$ & RMSE vs tolerance & 2026-09, green \\
GARCH(1,1) + EWMA & daily returns & constant mean & EWMA ($\lambda=0.94$) & rolling forecast vs realized & 2026-09, green \\
Black-76 + collar & F, K, vol, SOFR & European; discounted premiums & put-call parity; zero-cost gap & gap $\approx 0$ by construction & 2026-09, green \\
Participation note & F0, vol, SOFR, $p$ & risk-neutral martingale futures & MC vs closed form & MC re-check & 2026-09, green \\
IFRS 9 effectiveness & hedge P\&L pairs, ratios & principles-based (no 80--125\% band) & credit-dominance rejection & boundary tests & 2026-09, green \\
Discrete-CVaR ratio & scenario P\&L grid & discrete scenarios & MDPI (2025) hedge-ratio study & directional invariants & 2026-09, green \\
GARCH-BEKK cross hedge & coffee/gold returns & full 2$\times$2 BEKK(1,1), variance-targeted & literature hedge ratios & ratio stability & 2026-09, green \\
Margin stress & frozen path, vol, program & monetisable-put funding cap & FSB 2024 framing; relief identities & peaks vs buffer & 2026-09, green \\
Historical VaR/ES & realized P\&L & empirical tail & ES $\geq$ VaR invariant & rolling breach count & 2026-09, green \\
EE/EPE/PFE & F0, vol, SOFR, book & lognormal martingale & analytic forward MtM & invariants in suite & 2026-09, green \\
Unilateral CVA & EE profile, hazard, SOFR & flat hazard/SOFR; no DVA & QuantLib re-derivation & CVA-vs-QuantLib re-quote & 2026-09, green \\
\bottomrule
\end{tabular}
\caption{Model inventory (thirteen models, Phases 1--5; all reviewed 2026-09 with the suite green).}
\end{table}

\subsection{Conceptual soundness}

Each model's assumptions are stated in its module docstring and reproduced
in the inventory above. Two deliberate simplifications are flagged rather
than hidden: the Schwartz--Smith fit is reduced-form on a single snapshot
(volatility parameters would need the full Kalman MLE), and the collar's
margin treatment assumes the long put's gain is monetisable day-by-day
(real clearing charges full futures VM daily and settles options
separately).

\subsection{No-look-ahead gate}

Calibration windows are gated fail-closed: any observation after the
valuation as-of date raises a \texttt{LookaheadError}. The historical
as-of reconstructions in the outcomes program pre-truncate each frozen
series to the as-of window and then gate; an untruncated (contaminated)
input is refused. The gate is itself tested against a poisoned window.

\subsection{Outcomes analysis}

\begin{table}[h!]
\centering\small
\begin{tabular}{p{3.2cm}p{7.8cm}p{3.6cm}}
\toprule
\textbf{Check} & \textbf{Finding} & \textbf{Status} \\
\midrule
Working vol forecast & GARCH RMSE 4.01 vs EWMA 6.80 \%/yr over 4 rolling windows (ratio 0.59 $\leq$ 1.10) & PASS \\
No-look-ahead gate & active on 8 frozen series (as-of 2026-09-04); rejects a poisoned as-of window & PASS \\
Curve stability & front implied yield 14.8--40.8 \%/yr over 6 historical as-of dates; backwardated at 4 of 6 as-ofs (two historical contango episodes) & PASS \\
VaR coverage (95\%), static historical & 38 breaches vs 25.2 expected over 4 windows (ratio 1.51, band 0.5--1.5) & \textcolor{red}{FLAG} \\
VaR coverage (95\%), filtered EWMA & 37 breaches vs 33.7 expected over 673 days (ratio 1.10), remediating the static FLAG & PASS \\
Note MC vs closed form & $|$diff$| = 8.2\times10^{-5}$ per unit notional, s.e.\ $6.2\times10^{-4}$ ($\leq 2$ s.e.) & PASS \\
CVA vs QuantLib & relative difference $1.9\times10^{-16}$ (same discretization, tol $10^{-9}$) & PASS \\
\bottomrule
\end{tabular}
\caption{Outcome-analysis findings on the frozen data.}
\end{table}

\paragraph{The VaR coverage FLAG, and its remediation.} The static
historical 95\% VaR under-covers by roughly 50\% on this book: 38 breaches
against 25.2 expected over four 126-day windows. Two causes were
identified. First, the classic unconditional-historical-VaR failure mode
on fat-tailed, volatility-clustered returns: the trailing window's
empirical quantile understates the next window's tail. Second, a flaw in
the backtest itself: holding one stale VaR constant across a 126-day test
window conflates volatility-regime drift with model failure. The
remediation is implemented in the package (ticket 13-05,
\texttt{risk.filtered\_var}): a filtered estimator pricing each day with
the lagged EWMA variance ($\lambda = 0.94$) scaled by the lagged price.
Re-run on the same book it covers at 5.5\% of days versus the 5\%
expectation (37/673, ratio 1.10, inside the pre-declared band). The
static historical estimator is retained in the findings as the baseline
and its FLAG stands as the recorded evidence for why the filtered form is
the production estimator of this workbench; the acceptance band was not
moved.

\section{Decision: Hedged vs Unhedged Outcome}

Scenario engine: 50,000 lognormal martingale paths over one quarter
(0.25y), the same distribution Black-76 prices. Buyer's economics: the
roaster pays volume $\times f_T$ regardless; each hedge's terminal P\&L
reduces it. Note that the collar's strikes bound the \emph{hedge P\&L},
not the cost: between the strikes the cost is locked at the forward;
below the put (280) and above the call (382.06) the net cost moves 1:1
with the market again: a crash below the put keeps helping the buyer,
a spike above the call keeps hurting.

\begin{table}[h!]
\centering\small
\begin{tabular}{lrrrrr}
\toprule
\textbf{Strategy} & \textbf{Cost mean} & \textbf{Cost p95} & \textbf{Cost worst} &
\textbf{Margin peak mean} & \textbf{Margin peak p95} \\
\midrule
Unhedged & 3,895,696 & 5,288,885 & 10,296,484 & --- & --- \\
Futures & \textbf{3,891,000} & 3,891,000 & 3,891,000 & 43,655 & 103,797 \\
Collar & 3,893,639 & 4,595,138 & 9,602,737 & 30,647 & 44,250 \\
\bottomrule
\end{tabular}
\caption{Annual purchase cost (1.2M lb) and margin-funding peaks (USD),
2.67 average contracts, put 280 / call 382.06, 0.25y horizon.}
\end{table}

\paragraph{Reading the table.} The full futures hedge is a strict price
lock at \usd{}3,891,000 ($= 1.2\mathrm{M}\times 3.2425$); its residual
risk is purely margin funding: mean peak \usd{}43.7k, 95th percentile
\usd{}103.8k, comfortably inside the \usd{}800k reference-collateral
assumption carried from the Starbucks 10-K framing. The collar costs
almost the same on average (\usd{}3.894M) but trades tail exposure for
liquidity: its margin peak is capped at the put protection
($44{,}250 = (324.25 - 280)\times \tfrac{8}{3}\times 375$), roughly 58\% lower
at the 95th percentile, while the cost re-floats 1:1 beyond the strikes
(p95 \usd{}4.595M, worst \usd{}9.6M in-simulation). The unhedged roaster
owns the entire distribution: p95 \usd{}5.29M, simulated worst \usd{}10.3M.

\textbf{Decision.} For a roaster whose binding constraint is cash-flow
certainty (the Starbucks 10-K template: \usd{}37.9M of margin collateral
posted as of FY2025), the choice is between the futures lock (price
certainty, full margin exposure) and the collar (locked cost band,
half the margin peak, tail re-exposure). The workbench's numbers support
the collar for a margin-constrained buyer and the futures lock when
budget certainty dominates, with the IFRS~9 cash-flow-hedge
designation and effectiveness testing identical in structure for both.

\section{Limitations}\label{sec:limitations}

\begin{itemize}
\item \textbf{Basis risk is not modelled.} All hedges are evaluated on the
front KC=F series; the roaster's physical purchases settle on a differential
basis to the exchange. A basis-scenario overlay is the natural next step.
\item \textbf{Historical volatility only.} The working vol is a GARCH(1,1)
point estimate on frozen history; option premiums (collar legs, the note)
inherit its parameter risk. No implied-vol surface is available for free
coffee options.
\item \textbf{VaR remediation is EWMA-level.} The static historical VaR's
under-coverage is remediated by the filtered (EWMA $\lambda=0.94$) estimator
(Section~\ref{sec:validation}); a GARCH-conditional-vol scaling would be
the next refinement but is not implemented. The static estimator is kept
only as the flagged baseline.
\item \textbf{Unofficial source.} The Citi MQA Forage simulation is a
learning artifact, not a source of Citi methodology; all toy numbers from
it are labelled context only.
\item \textbf{Junior xVA only.} Unilateral CVA with flat hazard and flat
SOFR; no DVA/FVA/MVA, no wrong-way-risk modelling beyond a qualitative
discussion (a buyer short a call on coffee is long the reference asset;
WWR here is benign by inspection).
\item \textbf{Margin conventions are assumptions.} \usd{}8k initial margin
per contract and the monetisable-put funding cap are stated learner
assumptions, not exchange feeds.
\end{itemize}
\begin{thebibliography}{13}
\bibitem{black76} Black, F. (1976). The pricing of commodity contracts.
\emph{Journal of Financial Economics} 3(1--2), 167--179.
\bibitem{bollerslev86} Bollerslev, T. (1986). Generalized autoregressive
conditional heteroskedasticity. \emph{Journal of Econometrics} 31(3),
307--327.
\bibitem{sr262} Board of Governors of the Federal Reserve System, OCC, and
FDIC (2026). \emph{Supervisory Guidance on Model Risk Management},
SR~26-2 attachment, April 17, 2026.
\bibitem{ifrs9} International Accounting Standards Board (2014, as
amended). \emph{IFRS 9 Financial Instruments}, Section~6, Hedge
Accounting. IFRS Foundation.
\bibitem{emir3} European Union (2024). Regulation (EU) 2024/2987 (EMIR~3),
amending Regulation (EU) No~648/2012, in force 24 December 2024.
\bibitem{emir3rts} European Securities and Markets Authority (2026).
\emph{Final Report on draft RTS amending Commission Delegated Regulation
(EU) No~149/2013 to further detail the new EMIR clearing thresholds
regime}, ESMA74-1049116226-944, February 2026.
\bibitem{priips} European Union (2014). Regulation (EU) No~1286/2014 on
key information documents (PRIIPs).
\bibitem{riskmetrics96} RiskMetrics Group (1996). \emph{RiskMetrics
Technical Document}, 4th edition. J.P.~Morgan/Reuters, New York.
\bibitem{kupiec95} Kupiec, P.~H. (1995). Techniques for verifying the
accuracy of risk measurement models. \emph{Journal of Derivatives}
3(2), 73--84.
\bibitem{schwartzsmith00} Schwartz, E.~S., and Smith, J.~E. (2000).
Short-term variations and long-term dynamics in commodity prices.
\emph{Management Science} 46(7), 893--911.
\bibitem{hull} Hull, J.~C. \emph{Options, Futures, and Other Derivatives}.
Pearson (multiple editions).
\bibitem{quantlib} QuantLib contributors. \emph{QuantLib: a free/open-source
library for quantitative finance}. \url{https://www.quantlib.org}.
\bibitem{fredsofr} Federal Reserve Bank of St.\ Louis (2026). Secured
Overnight Financing Rate (SOFR), series SOFR. FRED.
\url{https://fred.stlouisfed.org/series/SOFR}.
\end{thebibliography}

\end{document}

